\documentclass[12pt,a4paper]{article}
\usepackage[UTF8]{ctex}
\usepackage{amsmath, amssymb}
\usepackage{geometry}
\usepackage{booktabs}
\usepackage{float}
\usepackage{enumitem}
\usepackage{hyperref}

\geometry{left=2.5cm, right=2.5cm, top=2.5cm, bottom=2.5cm}

\title{\textbf{模型假设（论文第 X 章素材）}\\[6pt]
\large 通用假设 + 四个问题分别的假设与合理性说明}
\author{2026 年第七届"华数杯"大学生数学建模竞赛 B 题}
\date{2026/08/10}

\begin{document}
\maketitle

% ============================================================
\section{模型假设}
% ============================================================

% ============================================================
\subsection{通用假设（适用于四个问题）}

\begin{enumerate}[leftmargin=2em]
    \item \textbf{硬模块假设}：所有模块为硬模块，几何尺寸固定、不可
    缩放或变形；允许的变换仅为整体旋转（问题一至三为 $90^\circ$
    旋转，问题四为 $0^\circ/90^\circ/180^\circ/270^\circ$ 四向旋转），
    旋转不改变模块面积。

    \item \textbf{不重叠与紧贴}：任意两模块不允许相互重叠，但允许
    共享边界（紧贴放置）；模块可放置于芯片内的任意整数坐标位置。

    \item \textbf{离散网格假设}：输入数据的模块尺寸、端子坐标均为
    整数，故所有模块坐标与轮廓尺寸取整数取值，不引入连续坐标的
    数值误差。

    \item \textbf{理想化物理环境}：不考虑布线拥塞、热分布、电源网格、
    时序、信号完整性等物理约束，模块间无任何空间间距要求（可零间距
    紧贴）。

    \item \textbf{数据可靠性}：竞赛提供的 \texttt{.blocks} /
    \texttt{.nets} / \texttt{.pl} 数据完整且不含噪声，直接采用，
    不做缺失或异常处理。

    \item \textbf{面积下界基准}：任意合法布局的包围盒面积不小于模块
    总面积 $\Sigma A$，该下界作为解质量的统一度量基准
    （面积利用率 $\eta = \Sigma A / (WH) \le 1$）。
\end{enumerate}

% ============================================================
\subsection{问题一假设（自由轮廓面积最小化）}

\begin{enumerate}[leftmargin=2em]
    \item \textbf{自由轮廓}：芯片轮廓不受约束，包围盒由模块摆放自然
    决定；
    \item \textbf{无连接关系}：模块之间不存在线网连接，布局仅由几何
    尺寸决定（\texttt{.nets}/\texttt{.pl} 不参与问题一目标）；
    \item \textbf{两级目标}：主目标为包围盒面积最小；面积相同时以
    长宽比更接近 1 为优（题目验证标准），二者以字典序处理。
\end{enumerate}

\textbf{合理性}：问题一为纯几何打包问题，自由轮廓假设与题面一致；
两级目标序关系与题目"在面积相同的情况下长宽比越接近 1 越好"的
描述完全对应。

% ============================================================
\subsection{问题二假设（固定轮廓 HPWL 最小化）}

\begin{enumerate}[leftmargin=2em]
    \item \textbf{固定正方形轮廓}：轮廓边长
    $\text{side} = \lceil\sqrt{\Sigma A\,(1+d)}\rceil$，死区比例
    $d = 0.15$，所有模块必须完全落入轮廓内；
    \item \textbf{线长一阶估计}：以半周长线长（HPWL）度量线长代价，
    与真实布线长度成正比，忽略绕线与拥塞对线长的非线性影响；
    \item \textbf{引脚中心}：块引脚取模块几何中心
    $p = (2x + w,\ 2y + h)$，终端引脚取坐标 $p = (2x,\ 2y)$；
    \item \textbf{端子固定}：\texttt{.pl} 中的端子位置固定，不参与
    布局决策，仅参与 HPWL 计算；
    \item \textbf{线网为超边}：每条线网连接多个引脚，引脚间距离
    以包围盒半周长近似（最大最小坐标之差求和）。
\end{enumerate}

\textbf{合理性}：HPWL 是布图规划领域最常用的线长代理指标；引脚中心
近似在模块尺寸远小于布线长度时误差可忽略；固定轮廓与题面公式一致。
若考虑布线拥塞，需引入网格密度惩罚，已列入模型推广方向。

% ============================================================
\subsection{问题三假设（最小可行死区比例 $d^*$）}

\begin{enumerate}[leftmargin=2em]
    \item \textbf{单调性}：死区比例 $d$ 增大 $\Rightarrow$ 轮廓边长
    $\text{side}(d)$ 单调不减 $\Rightarrow$ 可行域单调扩大，可行性
    判定 $d \mapsto \text{feasible}(d)$ 为单调布尔函数；
    \item \textbf{判定充分性}：启发式可行性判定存在假阴性（搜索不足
    误判不可行），假设\textbf{并行多种子判定}（任一可行即可行）可将
    假阴性概率降到可忽略水平；
    \item \textbf{搜索强度相关}：$d^*$ 是"在当前判定强度（种子数）下
    验证过的最小可行值"，论文表述以该限定为准；
    \item \textbf{上界可行}：$d = 0.15$（问题二基准）已知可行，故
    $d^* \in [0, 0.15]$。
\end{enumerate}

\textbf{合理性}：单调性由轮廓公式直接导出（$\lceil\cdot\rceil$ 单调）；
多种子判定将假阴性概率按 $P \approx \prod_{s=1}^{K} P_s$ 指数下降，
实测种子数 3$\to$15 使 $d^*$ 改善 23--45\%；假设 3 保证了结论的
严谨表述边界。

% ============================================================
\subsection{问题四假设（L/T 模块最小包围盒）}

\begin{enumerate}[leftmargin=2em]
    \item \textbf{矩形分解}：L/T 型模块可\textbf{精确}表示为固定相对
    位置的矩形并集（超块），分解矩形两两不重叠、间隙为零；
    \item \textbf{四向旋转}：模块整体可旋转 $0^\circ/90^\circ/180^\circ/
    270^\circ$，旋转后分解矩形相对位置不变（角点变换 + 归一化）；
    \item \textbf{矩形级合法性}：无重叠判定在分解矩形层面进行
    （超块内部矩形也参与检查）；
    \item \textbf{最优性范围}：穷举完备性仅在本题 4 模块实例下成立，
    其结论不直接外推至大规模实例（大规模沿用框架但为启发式）。
\end{enumerate}

\textbf{合理性}：题目附图的 T/L 模块均为一笔画多边形，可精确矩形
分解（凹角仅 90°）；四向旋转为整数网格下的完备旋转集合；穷举
86\,016 组合全部通过矩形级重叠检查，完备性与可达性共同支撑
全局最优证明。

% ============================================================
\subsection{假设的敏感性说明}

\begin{table}[H]
\centering
\scriptsize
\setlength{\tabcolsep}{3pt}
\caption{主要假设的影响与可放宽性}
\begin{tabular}{llll}
\toprule
假设 & 影响 & 敏感性 & 可放宽方向 \\
\midrule
硬模块 & 目标/约束定义基础 & 低（题面既定） & 软模块（长宽比连续可变） \\
忽略物理约束 & 仅几何优化 & 低--中 & 联合热/布线/网格约束（模型推广） \\
引脚中心近似 & HPWL 精度 & 低（模块远小于布线） & 精确引脚坐标 \\
判定充分性 & $d^*$ 数值 & \textbf{中--高}（种子相关） & 更高种子/贝叶斯选点（已列展望） \\
穷举完备性 & Q4 最优性 & 低（穷举覆盖） & 大规模实例退回启发式 \\
\bottomrule
\end{tabular}
\end{table}

\end{document}
